\lesson{speckit}{用 Spec Kit 串起文件与执行}
\diagram{p_speckit}
Spec 是需求规格；\href{https://github.com/github/spec-kit}{Spec Kit} 是组织规格、方案与任务的可选开源工具。先理解方法，再按需选用；小改动用短记录即可。
\topic{步骤与对应文件：以安装版本为准}
\begin{tabularx}{\linewidth}{@{}L{39mm}Y@{}}
\toprule 步骤 & 作用与需要检查的内容\\\midrule
\texttt{constitution} & 建立项目原则，保存在 \texttt{.specify/memory/constitution.md}\\
\texttt{specify / clarify} & 写清需求、澄清歧义，维护 \texttt{spec.md}\\
\texttt{plan} & 调研写入 \texttt{research.md}，方案写入 \texttt{plan.md}\\
\texttt{checklist} & 按需检查需求质量；通过不等于功能已实现\\
\texttt{tasks / analyze} & 生成 \texttt{tasks.md}；按需核对需求、方案、任务是否一致\\
\texttt{implement} & 按任务实施，运行检查，保留结果\\
\texttt{converge} & 实施后找遗漏、追加剩余任务；仍需实际验收\\
\bottomrule\end{tabularx}

\texttt{plan} 还生成 \texttt{data-model.md}（数据组织）、\texttt{quickstart.md}（验证步骤）；有对外接口时生成 \texttt{contracts/}（接口约定），供实现与验收使用。

当前默认技能入口：Codex 用 \texttt{\$speckit-specify}；Claude Code、Copilot 用 \texttt{/speckit-specify}。当前功能由 \texttt{.specify/}\allowbreak\texttt{feature.json} 或运行设置指定，切换 Git 分支不会自动切换它。
\begin{prompt}{按 Spec Kit 推进当前功能}
项目：[路径]，依据：[要求与记录]，本轮步骤：[步骤名]，允许修改：[范围]。核对安装版本、调用入口、当前功能和前置文件；未安装或前置缺失先给补齐方案，不擅自安装或跳步。\par
保留已有文件；澄清歧义，查证未知。需要自动测试时写入规格并拆成任务，每项对应需求与检查方法。\par
只有用户能调用时给出命令并等待执行结果。报告实际产物、检查、冲突及下一步；未调用不称已执行，要求矛盾不擅自取舍。
\end{prompt}
不熟悉命令时让 AI 查官方说明并核对入口。讨论后整理规格，人审核目标与验收，AI 按任务实施；需求变更同步方案与检查，避免新提示覆盖旧约定。文档齐全不等于完成。
